\chapter{Introdução}
\label{cap:introdução}

O envelhecimento da população é um dos desafios sociais mais prementes da atualidade. Em
Portugal, como na generalidade dos países europeus, a proporção de pessoas com mais de 65 anos
tem crescido de forma contínua, colocando uma pressão crescente sobre os sistemas de saúde e
de apoio social \parencite{ins2023envelhecimento}. Neste contexto, a tecnologia emerge como um
vetor de resposta: soluções de assistência digital permitem alargar a capacidade de
monitorização e acompanhamento de idosos sem substituir a presença humana dos cuidadores.

O presente relatório descreve o trabalho realizado no âmbito do estágio curricular na empresa
Kind Tech, uma \textit{startup} portuguesa que desenvolve o CompyAI, um assistente de voz
baseado em inteligência artificial para acompanhamento de pessoas idosas institucionalizadas.
A plataforma estabelece uma ponte entre o idoso, o cuidador e a administração da instituição,
oferecendo um canal de comunicação natural e acessível para o utente, e informação estruturada
sobre o seu bem-estar para os profissionais que o acompanham.

O estágio incidiu sobre o desenvolvimento da plataforma enquanto produto de software, com
foco na robustez da camada de serviços, na definição de uma estratégia de testes automatizados
e na construção de um processo de entrega em ambiente de produção. Estas três dimensões
constituem a espinha dorsal da engenharia de software do produto e determinam a sua
fiabilidade, segurança e capacidade de evolução sustentada.

%----------------------------------------------------------------------------------------
\section{Enquadramento e Contexto}
\label{sec:enquadramento}

A Kind Tech desenvolve o CompyAI como uma plataforma \textit{full-stack} com três camadas de
serviço independentes: uma interface de utilizador acessível a partir de qualquer dispositivo
com navegador web, um serviço central responsável pela lógica de negócio e pela persistência
de dados, e um microserviço dedicado ao processamento de inteligência artificial. Esta
separação de responsabilidades permite que cada camada evolua de forma autónoma, sem introduzir
dependências desnecessárias entre domínios distintos do sistema.

O fluxo principal da aplicação começa quando o idoso inicia uma conversa por voz. O áudio é
capturado no dispositivo e enviado para o serviço de inteligência artificial, que o transcreve
para texto, gera uma resposta em linguagem natural adequada ao perfil e ao historial do
utilizador, e converte essa resposta novamente em fala. Em paralelo, o serviço central persiste
a conversa e produz um sumário estruturado com os tópicos abordados, o tom emocional
detetado e os pontos-chave identificados, que fica disponível para consulta pelo cuidador.
O microserviço de inteligência artificial não está exposto diretamente à internet: apenas o
serviço central comunica com ele, numa rede interna isolada, o que reduz a superfície de
ataque do sistema.

A plataforma serve três perfis de utilizador com responsabilidades e interfaces distintas.
O \textbf{idoso} interage exclusivamente por voz, sem necessidade de literacia digital; o
assistente adapta o seu tom, vocabulário e duração das respostas ao perfil individual de cada
utilizador. O \textbf{cuidador} acede a um painel de monitorização onde visualiza o historial
de conversas, os sumários de bem-estar e os alertas gerados pelo sistema para os utentes sob
a sua responsabilidade. O \textbf{administrador} gere a instituição: cria e configura contas,
define parâmetros globais e tem uma visão transversal sobre toda a população da instituição.
Esta estrutura de papéis garante que cada utilizador acede apenas à informação e às ações
relevantes para a sua função.

No plano da infraestrutura, a plataforma é distribuída como um conjunto de serviços
contentorizados, alojados na nuvem, com acesso externo protegido por TLS. As atualizações
da aplicação e da base de dados são geridas de forma automatizada, minimizando a intervenção
manual e o risco de inconsistências entre versões do código e do esquema de dados.

No plano da qualidade de software, a estratégia de testes cobre diferentes níveis de
abstração: desde a validação da lógica de negócio isolada de dependências externas, até à
verificação dos contratos de dados partilhados entre serviços, passando pela validação das
integrações com fornecedores externos e pela confirmação de que as configurações críticas de
produção estão corretas antes de qualquer publicação.

O estágio insere-se, portanto, num produto em fase ativa de desenvolvimento, com utilizadores
reais e requisitos concretos de fiabilidade e segurança. O trabalho realizado contribui
diretamente para a maturidade da plataforma enquanto produto de software pronto para crescer.

%----------------------------------------------------------------------------------------
\section{Descrição do Problema}
\label{sec:problema}

As instituições de apoio a idosos debatem-se com um desequilíbrio estrutural entre a procura
de acompanhamento personalizado e os recursos humanos disponíveis. Um cuidador gere
frequentemente dezenas de utentes em simultâneo, tornando inviável uma interação contínua e
individualizada com cada um. Esta realidade traduz-se em isolamento social e em défice de
estimulação cognitiva para os idosos, com consequências documentadas no seu bem-estar físico
e mental \parencite{who2021ageism}.

As soluções tecnológicas existentes, aplicações de lembretes de medicação, sistemas de
alarme e plataformas de videochamada, não abordam o problema de forma integrada. Não
estabelecem um canal de comunicação natural e acessível para utilizadores com literacia
digital reduzida, não produzem informação estruturada e contínua sobre o estado emocional do
utente, e não oferecem ao cuidador uma visão longitudinal do bem-estar de cada idoso. A
lacuna não é apenas tecnológica; é também de desenho: as soluções existentes foram concebidas
para utilizadores com capacidades digitais que muitos idosos institucionalizados não possuem.

O CompyAI responde a este problema através de um assistente conversacional por voz — o canal
de comunicação mais natural e universal — que não exige qualquer interação com ecrãs ou
teclados. Em simultâneo, cada conversa é automaticamente analisada e transformada em
informação estruturada acessível ao cuidador, criando um registo longitudinal do bem-estar do
idoso sem impor trabalho adicional a nenhum dos intervenientes.

O desafio de engenharia de software consiste em construir uma plataforma que responda a três
requisitos em tensão: baixa latência nas interações de voz, para que a conversa seja fluida e
natural; integridade e confidencialidade dos dados clínicos, dado o carácter sensível da
informação processada; e fiabilidade de um serviço em produção contínua, onde uma
indisponibilidade tem impacto direto em utilizadores vulneráveis.

\subsection{Objetivos}
\label{subsec:objetivos}

O estágio tem como objetivo principal contribuir para a consolidação do CompyAI enquanto
produto de software robusto, seguro e evolutível. Este objetivo desdobra-se nos seguintes
objetivos específicos:

\begin{itemize}
    \item \textbf{Desenvolvimento do serviço central:} implementar e manter as
    funcionalidades da camada de serviços da plataforma, incluindo autenticação, gestão de
    perfis de utilizador, lógica de negócio associada às conversas e módulos de administração
    de instituições e contas.

    \item \textbf{Estratégia de testes automatizados:} definir e implementar uma abordagem de
    testes que cubra múltiplos níveis — lógica de negócio, contratos de dados, integrações
    externas e validação de configurações de produção — de forma a reduzir o risco de
    regressões e aumentar a confiança no processo de entrega.

    \item \textbf{Processo de entrega em produção:} conceber e operacionalizar a
    infraestrutura de produção na nuvem, incluindo a orquestração dos serviços, a segurança
    das comunicações e a automatização das atualizações do sistema.
\end{itemize}

\subsection{Abordagem}
\label{subsec:abordagem}

O desenvolvimento da plataforma segue uma metodologia ágil com iterações curtas, permitindo
incorporar rapidamente \textit{feedback} e ajustar prioridades face a novos requisitos
\parencite{beck2001agile}. Cada funcionalidade é desenvolvida de forma isolada, integrada na
base de código após revisão, e promovida a produção de forma controlada.

A arquitetura do serviço central foi desenhada com separação explícita entre a camada de
apresentação — que recebe e valida os pedidos externos —, a camada de aplicação — onde reside
a lógica de negócio —, e a camada de infraestrutura — que encapsula as dependências externas
como a base de dados e os serviços de terceiros \parencite{martin2017clean}. Esta organização
facilita a testabilidade de cada camada de forma independente e permite substituir
implementações concretas sem alterar a lógica de negócio.

A estratégia de testes segue a pirâmide proposta por \textcite{cohn2009succeeding},
privilegiando testes unitários rápidos na base e reservando os testes de integração para as
fronteiras do sistema, onde o risco de incompatibilidade com dependências externas é mais
elevado.

O processo de entrega baseia-se num modelo de infraestrutura declarativa, em que o estado
desejado do sistema de produção é descrito em ficheiros de configuração versionados, e não
configurado manualmente \parencite{humble2010continuous}. Isto garante que qualquer membro da
equipa pode recriar o ambiente de produção de forma determinística e que as atualizações são
reprodutíveis e auditáveis.

\subsection{Contributos}
\label{subsec:contributos}

O trabalho realizado no âmbito deste estágio traduz-se nos seguintes contributos para a
organização e para o produto:

\begin{itemize}
    \item Implementação do módulo de administração da plataforma, com gestão de instituições,
    contas de cuidadores e perfis de idosos, e controlo de acesso por perfil de utilizador.

    \item Definição e implementação de uma suite de testes automatizados que cobre os
    diferentes níveis da pirâmide de testes, aumentando a cobertura e reduzindo o risco de
    regressões em futuras iterações.

    \item Construção da infraestrutura de produção na nuvem, tornando o processo de entrega
    reprodutível, documentado e independente de intervenção manual.

    \item Automatização das atualizações do esquema de base de dados no arranque do sistema,
    eliminando um passo manual crítico no processo de atualização da plataforma.

    \item Contribuição para a maturidade do processo de entrega de software da equipa,
    estabelecendo práticas que a empresa poderá manter e evoluir de forma autónoma.
\end{itemize}

\subsection{Planeamento do Trabalho}
\label{subsec:planeamento}

O estágio decorreu entre outubro de 2025 e junho de 2026, com uma duração total de
aproximadamente 36 semanas, organizadas em quatro fases principais. O cronograma detalhado
encontra-se no Anexo~\ref{appendix:A}.

A \textbf{Fase 1} (semanas 1–6) correspondeu à integração na equipa, à familiarização com o
repositório e à compreensão da arquitetura e das práticas de desenvolvimento vigentes.
A \textbf{Fase 2} (semanas 7–18) centrou-se no desenvolvimento ativo de funcionalidades do
serviço central, nomeadamente os módulos de administração e os fluxos de integração de novos
utilizadores na plataforma.
A \textbf{Fase 3} (semanas 19–28) foi dedicada à estratégia de testes e à construção da
infraestrutura de produção.
A \textbf{Fase 4} (semanas 29–36) incidiu na estabilização da plataforma, na correção de
defeitos identificados em ambiente de produção e na elaboração do presente relatório.

%----------------------------------------------------------------------------------------
\section{Estrutura do Relatório}
\label{sec:estrutura}

Para além da presente introdução, este documento contém mais quatro capítulos.

No Capítulo~\ref{cap:estado_da_arte} é apresentada a revisão do estado da arte,
enquadrando o CompyAI no panorama das soluções de assistência digital a idosos e dos padrões
de arquitetura de software relevantes para o trabalho desenvolvido. São analisadas abordagens
existentes à assistência por voz em contexto de saúde, modelos de arquitetura para sistemas
com múltiplos serviços independentes, e as práticas de qualidade de software aplicáveis ao
contexto do projeto.

No Capítulo~\ref{cap:analise_desenho} é descrita a análise e o desenho da solução,
apresentando os requisitos funcionais e não funcionais da plataforma, a arquitetura do sistema
e as principais decisões de desenho que orientaram a implementação. É dado especial destaque
à organização dos serviços, ao modelo de dados e à estratégia de segurança adotada.

No Capítulo~\ref{cap:implementação} são apresentados os detalhes das funcionalidades
desenvolvidas durante o estágio: o módulo de administração da plataforma, a suite de testes
automatizados e a infraestrutura de entrega em produção. São discutidos os desafios encontrados
e as soluções adotadas em cada um destes domínios.

No Capítulo~\ref{cap:conclusão} são apresentadas as conclusões do trabalho, com uma avaliação
crítica dos resultados obtidos face aos objetivos definidos, uma reflexão sobre as
aprendizagens adquiridas e a identificação de trabalho futuro que poderá dar continuidade ao
que foi desenvolvido durante o estágio.
